\clearpage
\noindent\textbf{相遇抽取与收藏实现}\par\smallskip
\label{sec:impl-08-draw}
\par\medskip\noindent\textbf{待决定结果与加权选择。}
相遇模块将匹配候选转换成可决定的卡片。\DesignCode{draw()}先检查当前账号是否已有待决定结果；有则返回同一结果，不因刷新或重复点击重新抽取。没有待决定结果时才构造候选池，按匹配分形成累计票区间，抽中后保存快照。

当前版本每日抽取不限次数，但一次只能保留一个待决定结果，收藏上限为五位。初始加入的三位虚拟人物占用收藏位置；移除不会自动补回。没有可用候选或收藏已满时返回禁用原因，而不是制造一个不符合条件的人选。

\begin{table}[H]\centering\small
\caption{抽取和收藏的状态检查点}
\begin{tabularx}{\linewidth}{@{}p{29mm}p{25mm}X@{}}\toprule
当前条件 & 用户动作 & 处理结果 \\ \midrule
已有待决定卡片 & 再次抽取 & 复用原卡，不增添重复记录 \\
未达收藏上限 & 喜欢 & 保存收藏，决定记录只处理一次 \\
收藏已满 & 喜欢或抽取 & 拒绝越过容量约束 \\
不想继续展示 & 跳过或移除 & 进入冷却，不删除已发生的交流 \\ \bottomrule
\end{tabularx}\end{table}

\par\medskip\noindent\textbf{幂等与时间边界。}
\DesignCode{decide()}检查决定记录的账号归属和当前状态；同一请求重试不应额外占一个收藏位。\DesignCode{remove()}写入三十天冷却期，后续候选过滤仍需检查该状态。抽取日期使用UTC+8，与虚拟日期推进及服务端纸条日额度分别处理。

测试通过固定时间和可注入的随机票号检查选择边界、冷却、收藏满额、初始化仅一次和待决定复用。并发路径既依赖应用事务，也受唯一约束保护；不能只测单次成功抽取就声称并发一致性已通过。
